\VerdierPhysicalPageAnchor{0028}%
\VerdierSemanticAnchor{intro-7-locally-compact-model}
cas particulièrement simple des espaces localement compacts fournit le
modèle des énoncés qu'on obtient. Soient $X$ et $Y$ deux espaces localement
compacts (disons : de dimension cohomologique finie), $A$ un anneau
(noethérien), $\mathbf D(X_A)$ et $\mathbf D(Y_A)$ les catégories dérivées
des catégories de faisceaux de $A$-modules $X_A$ et $Y_A$ sur $X$ et $Y$
respectivement, $f\colon X\longrightarrow Y$ une application continue. On
peut alors définir dans cette situation :

\VerdierSemanticAnchor{intro-7-f-shriek}
\begin{enumerate}
\renewcommand{\labelenumi}{\alph{enumi})}
\item Un foncteur
\[
f_!\colon X_A\longrightarrow Y_A,
\]
image directe à support propre, et le foncteur dérivé total correspondant :
\VerdierSemanticAnchor{intro-7-rf-shriek}
\[
\mathrm Rf_!\colon \mathbf D(X_A)\longrightarrow \mathbf D(Y_A).
\]
(Lorsque $f\colon X\longrightarrow Y$ est une application propre, le
foncteur $f_!$ est le foncteur image directe usuel.)

\VerdierSemanticAnchor{intro-7-f-upper-shriek}
\item Un foncteur
\[
f^!\colon \mathbf D^+(Y_A)\longrightarrow \mathbf D^+(X_A)
\]
(le signe $+$ signifie qu'on se limite aux complexes bornés inférieurement)
dont la définition précise est assez technique, mais qui est caractérisé, à
isomorphisme unique près, par la propriété c).

\VerdierSemanticAnchor{intro-7-duality-adjunction}
\item Un isomorphisme bifonctoriel :
\VerdierSemanticAnchor{intro-7-delta-f}
\[
\Delta_f(F,G)\colon
\operatorname{RHom}(\mathrm Rf_!F,G)
\longrightarrow
\operatorname{RHom}(F,f^!G)
\]
défini pour tout objet $F$ de $\mathbf D(X_A)$ et tout objet $G$ de
$\mathbf D^+(Y_A)$ ($\operatorname{RHom}$ désigne le foncteur dérivé total
du foncteur \og complexe des homomorphismes\fg{}).
\end{enumerate}

Les phénomènes importants sont alors les suivants :

\VerdierSemanticAnchor{intro-7-duality-D1}
\noindent\textit{D1)} La formation du foncteur $f^!$ est transitive pour $f$
variable.

\VerdierSemanticAnchor{intro-7-duality-D2}
\noindent\textit{D2)} L'isomorphisme $\Delta_f$ est compatible avec la
composition des applications.

\VerdierSemanticAnchor{intro-7-duality-D3}
\noindent\textit{D3)} Lorsque $X$ est lisse sur $Y$ de dimension relative
$n$, \emph{i.e.} lorsque, localement sur $X$, $X$ est $Y$-isomorphe au
produit d'un ouvert de $Y$ par l'espace topologique $\mathbf R^n$ (exemple :
$Y$ = un point, $X$ variété topologique), on a pour tout objet $G$ de
$\mathbf D^+(Y_A)$ un isomorphisme canonique :
\VerdierSemanticAnchor{intro-7-relative-orientation}
\[
f^!(G)\simeq f^*(G)\otimes\omega_{X/Y}[n],
\]
où $f^*(G)$ est l'image inverse usuelle, $\omega_{X/Y}$ un faisceau
localement libre de rang $1$ : le faisceau d'orientation relative de $X$ sur
$Y$, et où $\omega_{X/Y}[n]$ désigne le complexe obtenu en faisant subir au
complexe $\omega_{X/Y}$ un décalage des degrés de $n$ unités.
